\section{\texorpdfstring{$\mathbb{Z}$}{Z} Versus \texorpdfstring{$\mathbb{Z}/m\mathbb{Z}$}{Z/mZ}}

As definições de grupo e subgrupo tornam claro que a interseção de qualquer coleção --- 
ainda que infinita --- de subgrupos de um grupo $G$ é também um subgrupo de $G$. (De fato, 
o elemento identidade pertence à interseção, e se $g_1$ e $g_2$ pertencem à interseção, então 
$g_1 g_2 \text{ e } g_1^{-1}$ também pertence).

\begin{definicao}
Sejam $a, b, \ldots, c$ elementos de um grupo $G$. Dizemos que a interseção $H$ de todos os 
subgrupos de $G$ (incluindo o próprio $G$) que contêm $a, b, \ldots, c$ é o \textbf{subgrupo gerado} por 
$a, b, \ldots, c$ e que estes formam um conjunto gerador de $H$. Neste caso, escreveremos
$H = \langle a, b, \ldots, c \rangle$.
\end{definicao}

Isto pode ser expresso ao dizermos que $\langle a, b, \ldots, c \rangle$ é o menor subgrupo de $G$ 
contendo $a, b, \ldots, c$; pode acontecer dele ser o próprio $G$ (e.g., no grupo aditivo $\mathbb{Z}$, 
temos $\langle 2, 3 \rangle = \mathbb{Z}$).

Seja $G$ um grupo e seja $x$ um elemento de $G$. Se $G$ é um grupo aditivo, escrevemos $0x$ para o elemento 
neutro $0$ de $G$, $-x$ para o elemento $x'$ definido por $x' + x = 0$, escrevemos $nx$, quando $n$ é um inteiro 
$> 0$, para a soma $x + \cdots + x$ de $n$ termos iguais a $x$, e $(-n)x$ para $-(nx)$. Por indução vemos que
$\{ nx : n \in \mathbb{Z} \}$ está contido em $\langle x \rangle$. Também por indução verifica-se as fórmulas
\[
mx + nx = (m+n)x
\quad \text{e} \quad
m(nx) = (mn)x,
\]
para quaisquer $m, n \in \mathbb{Z}$. A primeira fórmula diz que $\{ nx : n \in \mathbb{Z} \}$ é um subgrupo de $G$. Logo,
$\{ nx : n \in \mathbb{Z} \}$ contém $\langle x \rangle$, e portanto deve ser igual a $\langle x \rangle$.
Por conveniência, enunciamos isto como uma proposição apenas no caso em que $G$ é multiplicativo. Neste caso, 
escrevemos $x^0$ para o elemento neutro $1$ de $G$, $x^{-1}$ para o elemento $x'$ de $G$ definido por $x x' = 1$, 
escrevemos $x^n$, quando $n$ é um inteiro $> 0$, para o produto
$x \cdots x$ de $n$ fatores iguais a $x$, e $x^{-n}$ para $(x^n)^{-1}$

\begin{proposicao}
Se $G$ for um grupo multiplicativo, então para cada elemento $x$ de $G$, o subgrupo gerado por $x$ consiste 
em todos os elementos $x^n$ com $n \in \mathbb{Z}$.
\end{proposicao}

Sejam $G$ como na Proposição 11.2, e seja $M$ o conjunto de todos os inteiros $a$ tais que $x^a = 1$. Como 
$x^0 = 1$, $M$ é não-vazio. Além disso, para qualquer inteiro $a$ e qualquer inteiro $b$, a fórmula
$x^{a-b} = x^a (x^b)^{-1}$ mostra que $x^a = x^b$ se, e somente se, $a - b \in M$. Em particular, $M$ é 
fechado para a subtração. Portanto, $M$ satisfaz as hipóteses da Proposição 3.2 (em outras palavras, ele 
é um subgrupo de $\mathbb{Z}$) e consiste em todos os múltiplos de algum inteiro $m \ge 0$; se $m \ne 0$, 
ele é o menor inteiro $> 0$ tal que $x^m = 1$. Logo, se $m = 0$, então todos os elementos $x^a$ são distintos; 
e se $m > 0$, $x^a$ é igual a $x^b$ se, e somente se, $a \equiv b \pmod{m}$.

\begin{definicao}
Um isomorfismo entre dois grupos $G$ e $G'$ é uma aplicação bijetiva
$f : G \to G'$
tal que $f$ transforma a operação $\ast_G$ do grupo $G$ na operação $\ast_{G'}$ do grupo $G'$(\emph{i.e.},
$f(x \ast_G y) = f(x) \ast_G' f(y) \text{ para quaisquer } x, y \in G$). 
\end{definicao}

Por exemplo, entre os grupos aditivos $\mathbb{Z}/2014\mathbb{Z}$ e $\mathbb{Z}/2015\mathbb{Z}$ não existe 
isomorfismo algum. Por outro lado, temos o seguinte isomorfismo $f$ entre o grupo multiplicativo 
$(\mathbb{Z}/9\mathbb{Z})^\times$ e o grupo aditivo $\mathbb{Z}/6\mathbb{Z}$.

\begin{center}
$f((2^n \text{ mod } 9)) = (n \text{ mod } 6)$.    
\end{center}

De fato, cada elemento de $(\mathbb{Z}/9\mathbb{Z})^\times$ é da forma $(2^n \bmod 9)$, a aplicação é bijetiva, e
$f\big((2^{n_1} \bmod 9)\big) + f\big((2^{n_2} \bmod 9)\big)
= (n_1 + n_2 \bmod 6) = f\big((2^{n_1+n_2} \bmod 9)\big).$

O conceito de isomorfismo pode ser transportado de uma maneira óbvia para anéis e corpos.

\begin{proposicao}
Seja $G$ um grupo multiplicativo gerado por um elemento $x$. Então: ou $G$ é infinito e a aplicação
$x^n \mapsto n$ é um isomorfismo entre $G$ e o grupo aditivo $\mathbb{Z}$, ou $G$ é finito consistindo 
em $m$ elementos e a aplicação $x^n \mapsto (n \bmod m)$ é um isomorfismo entre $G$ e o grupo aditivo 
$\mathbb{Z}/m\mathbb{Z}$.
\end{proposicao}

Claramente, se $G$ é um grupo qualquer e $x$ é um elemento qualquer de $G$, a Proposição 11.4 pode ser aplicada 
para o subgrupo de $G$ gerado por $x$.

\begin{definicao}
O número de elementos de um grupo finito é chamado de sua \underline{ordem}. Se um grupo finito puder ser gerado por 
um elemento, dizemos que ele é \underline{cíclico}. Se um elemento $x$ de um grupo qualquer gera um grupo finito de 
ordem $m$, dizemos que $m$ é a ordem de $x$.
\end{definicao}

Por exemplo, a ordem do grupo aditivo $\mathbb{Z}/m\mathbb{Z}$ é $m$ e a ordem do grupo multiplicativo 
$(\mathbb{Z}/m\mathbb{Z})^\times$ é $\varphi(m)$. O grupo $\mathbb{Z}/m\mathbb{Z}$ é cíclico gerado por $(1 \bmod m)$ 
ou $(-1 \bmod m)$, enquanto que o grupo $(\mathbb{Z}/8\mathbb{Z})^\times$ possui ordem $4$ e não é cíclico. De fato, 
este último grupo não possui elemento de ordem $4$.

\[
\begin{array}{c|cccc}
\text{elemento em } (\mathbb{Z}/8\mathbb{Z})^\times 
& (1 \bmod 8) & (3 \bmod 8) & (5 \bmod 8) & (7 \bmod 8) \\ \hline
\text{ordem} & 1 & 2 & 2 & 2
\end{array}
\]

(Nota-se, contudo, que qualquer subconjunto de dois elementos de $(\mathbb{Z}/8\mathbb{Z})^\times$ que não contenha 
$(1 \bmod m)$ gera todo o grupo.)
